\begin{titlepage}
    \section*{Abstract}
        % Set the Scene
        Low-cost single-board computers like the Raspberry Pi series, especially 
        Raspberry Pi 4 and 5, are compact, energy-efficient platforms designed 
        to deliver versatile computing across a wide range of applications. They 
        rely on main memory to provide fast, temporary storage for actively processed 
        data. Each bit of data in \gls{dram} is stored in a single memory cell that 
        represents a binary value through the presence or absence of electrical 
        charge in capacitor~\cite{paperFlippingBits, paperRowhammeronaRPi, paperDrammer}.
        % Conflict
        As \gls{dram} cell dimensions continue to shrink in order to increase memory 
        density within limited physical space, technological progress enhances 
        capacity and efficiency while simultaneously amplifying susceptibility 
        to disturbance effects. \gls{dram} must be able to operate reliably without 
        compromising integrity due to charge loss. Repeated gate-induced 
        drain- and subthreshold leakage~\cite{paperFlippingBits, paperRHinDRAM} can cause sufficient charge 
        loss and disturbance, ultimately leading memory cells to flip 
        from 0 to 1 or vice versa. Such Software-induced fault, namely as 
        Rowhammer, deliberately exploit these effects by triggering accelerated 
        charge loss and disturbance errors, thereby corrupting memory contents 
        in a way that can enable privilege escalation and other severe security 
        consequences.
        % Experiments
        In this paper, we systematically evaluate the feasibility of mounting 
        rowhammer on modern Raspberry Pi platforms. To enable uncached accesses 
        to \gls{dram}, we examined multiple cache bypass techniques, of which cache 
        eviction strategies and ARM cache maintenance instructions proved effective. 
        Furthermore, we performed targeted accesses to specific \gls{dram} banks and 
        rows by reverse engineering the \gls{dram} address mapping using the DRAMA 
        approach and subsequently verifying the derived mappings.
        % Climax Experiment
        To outrun \gls{dram} refresh and circumvent \gls{trr}-based mitigation's, we ported 
        and adapted TRRespass to identify efficient hammering access patterns 
        capable of inducing bit flips to test the vulnerability.
        % Validate Climax Experiment
        Under the evaluated conditions, we could not observe any susceptibility 
        on modern Raspberry Pi models varying in model and ram configuration to 
        Rowhammer. However, factors such as potential \gls{trr} behavior and the 
        influence of \gls{dram} page policies preclude a definitive claim of immunity.
        % Return the protagonist to initial scene
        The results show that \gls{dram} in modern Raspberry Pi models remains intact 
        under tested Rowhammer conditions, demonstrating no observable compromise 
        of memory integrity.
        \\
        \textit{\textbf{Keywords}: Rowhammer, \gls{dram}, Raspberry Pi, Address Mapping, \gls{trr}}

\end{titlepage}
